% Manual de diseño de la plantilla UTI para trabajos de titulación
% Compilar desde la raíz del proyecto con:
%   latexmk -pdf -outdir=manual manual/Manual_de_Diseno.tex
\documentclass[12pt,a4paper,oneside]{report}

\usepackage[T1]{fontenc}
\usepackage[utf8]{inputenc}
\usepackage[spanish,es-tabla]{babel}
\usepackage{lmodern}
\usepackage{geometry}
\geometry{top=3cm,bottom=3cm,left=4cm,right=3cm}
\usepackage{setspace}
\onehalfspacing
\usepackage{array}
\usepackage{longtable}
\usepackage{tabularx}
\usepackage{xcolor}
\usepackage{graphicx}
\usepackage{listings}
\usepackage{enumitem}
\usepackage{amssymb}
\usepackage{fancyhdr}
\usepackage[hidelinks]{hyperref}

\definecolor{UTIMorado}{HTML}{422568}
\definecolor{UTINaranja}{HTML}{F56600}
\definecolor{CodigoFondo}{HTML}{F4F2F7}

\hypersetup{
  pdftitle={Manual de Diseño de la Plantilla para Trabajos de Integración Curricular},
  pdfauthor={Carrera de Ingeniería en Tecnologías de la Información},
  colorlinks=true,
  linkcolor=UTIMorado,
  urlcolor=UTIMorado,
  citecolor=UTIMorado
}

\pagestyle{fancy}
\fancyhf{}
\fancyhead[L]{\small Manual de Diseño}
\fancyhead[R]{\small Ingeniería en Tecnologías de la Información}
\fancyfoot[C]{\thepage}
\renewcommand{\headrulewidth}{0.4pt}
\renewcommand{\footrulewidth}{0pt}
\setlength{\headheight}{14pt}

\fancypagestyle{plain}{%
  \fancyhf{}%
  \fancyfoot[C]{\thepage}%
  \renewcommand{\headrulewidth}{0pt}%
}

\setlength{\parindent}{1.25cm}
\setlength{\parskip}{0pt}
\setlength{\emergencystretch}{2em}
\setlist{itemsep=0.25em,topsep=0.4em}

\lstset{
  basicstyle=\ttfamily\small,
  backgroundcolor=\color{CodigoFondo},
  frame=single,
  rulecolor=\color{UTIMorado},
  breaklines=true,
  columns=fullflexible,
  showstringspaces=false,
  keepspaces=true,
  upquote=true,
  literate=
    {á}{{\'a}}1 {é}{{\'e}}1 {í}{{\'i}}1 {ó}{{\'o}}1 {ú}{{\'u}}1
    {Á}{{\'A}}1 {É}{{\'E}}1 {Í}{{\'I}}1 {Ó}{{\'O}}1 {Ú}{{\'U}}1
    {ñ}{{\~n}}1 {Ñ}{{\~N}}1 {ü}{{\"u}}1 {¿}{{?`}}1
}

\newcommand{\criterio}[1]{\noindent\textbf{Criterio:} #1\par}
\newcommand{\resultado}[1]{\noindent\textbf{Resultado esperado:} #1\par}

\begin{document}
\hypersetup{pageanchor=false}

\begin{titlepage}
  \thispagestyle{empty}
  \begin{center}
    \vspace*{1.5cm}
    {\color{UTIMorado}\rule{\textwidth}{2pt}}\\[1.2cm]
    {\Large\bfseries UNIVERSIDAD INDOAMÉRICA\par}
    \vspace{0.5cm}
    {\large\bfseries FACULTAD DE INGENIERÍAS\par}
    \vspace{0.3cm}
    {\large\bfseries CARRERA DE INGENIERÍA EN TECNOLOGÍAS DE LA INFORMACIÓN\par}
    \vfill
    {\Huge\bfseries\color{UTIMorado} MANUAL DE DISEÑO\par}
    \vspace{0.5cm}
    {\LARGE Plantilla para Trabajos de Integración Curricular\par}
    \vspace{1cm}
    {\large Documento técnico para revisión y aprobación del Comité Curricular\par}
    \vfill
    \begin{tabular}{rl}
      \textbf{Código documental:} & UTI--TIC--MD--01 \\
      \textbf{Versión:} & 1.0 \\
      \textbf{Estado:} & Propuesta para aprobación \\
      \textbf{Fecha:} & 20 de julio de 2026 \\
    \end{tabular}
    \vfill
    {\color{UTINaranja}\rule{\textwidth}{2pt}}
  \end{center}
\end{titlepage}

\hypersetup{pageanchor=true}
\pagenumbering{roman}

\chapter*{Control del documento}
\addcontentsline{toc}{chapter}{Control del documento}

\begin{longtable}{|p{3.8cm}|p{8.8cm}|}
  \hline
  \textbf{Elemento} & \textbf{Definición} \\
  \hline
  Propietario del formato & Comité Curricular de la Carrera de Ingeniería en Tecnologías de la Información. \\
  \hline
  Responsable técnico & Coordinación de titulación o la unidad designada por el Comité Curricular. \\
  \hline
  Usuarios & Estudiantes, docentes tutores, lectores, miembros de tribunal y personal de apoyo académico. \\
  \hline
  Documento controlado & \texttt{uti\_tfm.sty}, estructura de carpetas y archivos base incluidos en el proyecto. \\
  \hline
  Finalidad & Normalizar la presentación de los Trabajos de Integración Curricular de la carrera. \\
  \hline
  Estado de esta versión & Propuesta técnica sujeta a revisión, observaciones y aprobación del Comité Curricular. \\
  \hline
\end{longtable}

\tableofcontents
\clearpage
\pagenumbering{arabic}

\chapter{Objeto, alcance y autoridad del manual}

\section{Objeto}

Este manual establece las decisiones de diseño, la arquitectura documental, las reglas tipográficas, los componentes reutilizables y los criterios de aceptación de la plantilla LaTeX para Trabajos de Integración Curricular de la Carrera de Ingeniería en Tecnologías de la Información.

El documento sirve como referencia técnica para que el Comité Curricular pueda evaluar la coherencia del formato, aprobar una versión institucional y controlar futuras modificaciones.

\section{Alcance}

El manual cubre:

\begin{itemize}
  \item configuración de página y tipografías;
  \item portada y páginas preliminares;
  \item estructura sugerida de capítulos;
  \item tablas, figuras, imágenes, ecuaciones, índices, anexos y referencias;
  \item metadatos para uno o dos autores;
  \item selección de líneas y sublíneas de investigación;
  \item compilación, validación, control de cambios y pruebas de aceptación.
\end{itemize}

No determina el contenido científico específico de cada proyecto ni reemplaza la evaluación académica del tutor, lectores o tribunal.

\section{Jerarquía normativa}

Cuando exista diferencia entre este manual y una resolución, instructivo o disposición institucional posterior, prevalecerá el documento institucional vigente. Toda modificación aprobada deberá reflejarse simultáneamente en el archivo de estilo, la plantilla, el Manual de Diseño y el Manual de Uso.

\chapter{Arquitectura del proyecto}

\section{Estructura controlada}

La plantilla separa la presentación, los metadatos y el contenido. La estructura base es la siguiente:

\begin{lstlisting}
main.tex
uti_tfm.sty
referencias.bib
Logos/
  LOGO-UTI-1.png
Preliminares/
  01_RepositorioDigital.tex
  02_AprobacionTutor.tex
  03_DeclaracionAutenticidad.tex
  04_AprobacionLectores.tex
  05_Dedicatoria.tex
  06_Agradecimiento.tex
  07_Indices.tex
  08_Resumen.tex
  09_Abstract.tex
Capitulos/
  01_Introduccion.tex
  02_Metodologia.tex
  03_Implementacion.tex
  04_Conclusiones.tex
  05_Referencias.tex
  06_Glosario.tex
  fig/
  img/
Anexos/
  01_Anexo.tex
manual/
  Manual_de_Diseno.tex
  Manual_de_Uso.tex
\end{lstlisting}

Los nombres reales de varias carpetas contienen tildes. Deben respetarse exactamente, incluidas mayúsculas, tildes y extensión del archivo.

\section{Responsabilidad de cada componente}

\begin{longtable}{|p{3.6cm}|p{9cm}|}
  \hline
  \textbf{Componente} & \textbf{Responsabilidad} \\
  \hline
  \texttt{main.tex} & Define metadatos, selecciona la línea de investigación y organiza el orden de inclusión de todas las partes. \\
  \hline
  \texttt{uti\_tfm.sty} & Centraliza márgenes, fuentes, títulos, numeración, portada, páginas especiales y comandos institucionales. \\
  \hline
  \texttt{Preliminares/} & Contiene documentos legales, dedicatorias, agradecimientos, índices, resumen y abstract. \\
  \hline
  \texttt{Capítulos/} & Contiene el desarrollo académico y técnico del trabajo. \\
  \hline
  \texttt{Anexos/} & Contiene material complementario registrado en el índice de anexos. \\
  \hline
  \texttt{referencias.bib} & Base bibliográfica utilizada por BibTeX con el estilo IEEEtran. \\
  \hline
  \texttt{Logos/} & Almacena los recursos gráficos institucionales autorizados. \\
  \hline
\end{longtable}

\chapter{Especificación visual}

\section{Página y márgenes}

\begin{longtable}{|p{4cm}|p{4cm}|p{4.4cm}|}
  \hline
  \textbf{Propiedad} & \textbf{Valor aprobado en la plantilla} & \textbf{Implementación} \\
  \hline
  Tamaño de papel & A4 & Clase \texttt{report}, tamaño 12 puntos, A4 y una cara. \\
  \hline
  Margen superior & 3 cm & Paquete \texttt{geometry}. \\
  \hline
  Margen inferior & 3 cm & Paquete \texttt{geometry}. \\
  \hline
  Margen izquierdo & 4 cm & Permite encuadernación y lectura institucional. \\
  \hline
  Margen derecho & 3 cm & Mantiene el área útil definida. \\
  \hline
  Encabezado & Sin contenido visible & Se elimina la regla del encabezado. \\
  \hline
  Pie de página & Número centrado & Aplicado también al estilo \texttt{plain}. \\
  \hline
\end{longtable}

\section{Tipografía y párrafo}

\begin{itemize}
  \item El cuerpo utiliza Century Gothic cuando se compila con XeLaTeX o LuaLaTeX y la fuente está instalada.
  \item Con pdfLaTeX se utiliza TeX Gyre Adventor como alternativa compatible.
  \item La portada, el resumen y el abstract utilizan Times New Roman; cuando no está disponible se emplea TeX Gyre Termes o NewTX como sustituto métrico.
  \item El tamaño base del documento es de 12 puntos.
  \item El cuerpo se configura con interlineado de 1,5; la portada, el resumen y el abstract usan interlineado sencillo.
  \item Los párrafos se justifican, usan sangría inicial de 1,25 cm y no añaden espacio vertical adicional entre párrafos.
\end{itemize}

\section{Capítulos y apartados}

Los capítulos se numeran con romanos y se presentan como \textbf{CAPÍTULO I}, \textbf{CAPÍTULO II}, etc. El título del capítulo aparece centrado y en mayúsculas. Los capítulos no numerados, como referencias y glosario, no muestran la palabra ``CAPÍTULO''.

Las secciones muestran numeración arábiga vinculada al capítulo: \textbf{1.1}, \textbf{1.2}, \textbf{2.1}, etc. Las subsecciones emplean un tercer nivel, por ejemplo \textbf{1.1.1}. Los encabezados de capítulo conservan números romanos y el Índice de Contenidos los identifica explícitamente como \textbf{CAPÍTULO I}, \textbf{CAPÍTULO II}, etc.

\section{Numeración de páginas}

\begin{enumerate}
  \item La portada no muestra número ni cuenta en la numeración.
  \item Las páginas preliminares comienzan en romano minúsculo mediante el comando institucional \texttt{UTIFrontMatter}.
  \item El cuerpo cambia a números arábigos mediante \texttt{UTIMainMatter}, sin reiniciar el contador.
  \item La variante \texttt{UTIMainMatterReset} reinicia el cuerpo en 1 y solo debe habilitarse por decisión institucional.
\end{enumerate}

\chapter{Diseño de la portada}

La portada se genera con \texttt{\textbackslash UTIPortada} y toma nueve parámetros. Su composición responde a los siguientes criterios:

\begin{longtable}{|p{4.2cm}|p{8.4cm}|}
  \hline
  \textbf{Elemento} & \textbf{Especificación} \\
  \hline
  Logo & Centrado, con caja máxima de 3 cm por 3 cm y conservación de proporciones. \\
  \hline
  Universidad & ``UNIVERSIDAD INDOAMÉRICA'', negrita, 16 puntos. \\
  \hline
  Facultad & Mayúsculas, negrita, 14 puntos. \\
  \hline
  Carrera & Mayúsculas, negrita, 12 puntos. \\
  \hline
  Tema & Etiqueta en negrita, línea horizontal de 1 punto y título en mayúsculas. \\
  \hline
  Modalidad & Texto de la modalidad en 12 puntos. \\
  \hline
  Autor y tutor & Bloque ubicado en la mitad derecha del área de texto, adaptable a uno o dos autores. \\
  \hline
  Ciudad y año & Centrados en la parte inferior; ciudad y país se separan mediante raya. \\
  \hline
\end{longtable}

La portada se encapsula en \texttt{titlepage}, no presenta pie de página y establece el contador en cero al finalizar.

\chapter{Páginas preliminares}

\section{Orden establecido}

\begin{enumerate}
  \item Autorización para consulta, reproducción y publicación electrónica.
  \item Aprobación del tutor o tutora.
  \item Declaración de autenticidad.
  \item Aprobación de lectores.
  \item Dedicatoria.
  \item Agradecimiento.
  \item Índices.
  \item Resumen ejecutivo.
  \item Abstract.
\end{enumerate}

Los textos legales reutilizan los metadatos definidos en \texttt{main.tex}. El estilo ajusta pronombres, etiquetas y conjugaciones cuando existe uno o dos autores y según los valores \texttt{masculino} o \texttt{femenino}.

\section{Índices}

La plantilla genera índices independientes de contenidos, tablas, figuras, imágenes, ecuaciones y anexos. Para que aparezcan correctamente deben realizarse compilaciones sucesivas hasta que los archivos auxiliares se estabilicen.

\section{Resumen y abstract}

Ambas páginas utilizan una composición equivalente y la fuente de portada:

\begin{itemize}
  \item nombre de la Universidad sin subrayado;
  \item facultad y carrera;
  \item título del proyecto, autor o autores y tutor;
  \item título central de la sección;
  \item un solo párrafo de hasta 300 palabras;
  \item entre tres y cuatro descriptores o palabras clave en orden alfabético.
\end{itemize}

El abstract cambia automáticamente las etiquetas a \textbf{TOPIC}, \textbf{AUTHOR}, \textbf{ADVISOR} y \textbf{KEYWORDS}; además usa el título inglés parametrizado.

\chapter{Estructura académica del cuerpo}

\begin{longtable}{|p{3cm}|p{4.2cm}|p{5.4cm}|}
  \hline
  \textbf{Capítulo} & \textbf{Título} & \textbf{Contenido esperado} \\
  \hline
  I & Introducción & Contextualización, estado del arte, problema, justificación, objetivos, alcance y fundamentación teórica. \\
  \hline
  II & Metodología & Diagnóstico, área de estudio, metodología, recolección de información, requerimientos, diseño y justificación técnica. \\
  \hline
  III & Implementación, Pruebas y Resultados & Solución, proceso de implementación, pruebas, validación, resultados, cronograma y presupuesto. \\
  \hline
  IV & Conclusiones y Recomendaciones & Conclusiones, recomendaciones y trabajos futuros. \\
  \hline
  No numerado & Referencias & Fuentes citadas mediante IEEEtran y BibTeX. \\
  \hline
  No numerado & Glosario & Definiciones de términos y siglas relevantes. \\
  \hline
  Anexos & Anexo A, B, C, etc. & Evidencias, instrumentos, configuraciones y documentación complementaria. \\
  \hline
\end{longtable}

Los textos entre corchetes y en cursiva son instrucciones de plantilla. Deben ser reemplazados por contenido definitivo antes de la entrega.

\chapter{Metadatos y comandos públicos}

\section{Datos centrales}

\begin{longtable}{|p{5.2cm}|p{7.4cm}|}
  \hline
  \textbf{Comando} & \textbf{Función} \\
  \hline
  {\scriptsize\ttfamily\textbackslash UTIDatosTituloIngles\{...\}} & Registra el título utilizado en el abstract. Si se omite, el sistema utiliza el título español. \\
  \hline
  {\scriptsize\ttfamily\textbackslash UTISeleccionInvestigacion\{L\}\{S\}} & Selecciona una línea y una sublínea del catálogo institucional. \\
  \hline
  {\scriptsize\ttfamily\textbackslash UTIDatosAutores} & Registra uno o dos autores y su género gramatical. \\
  \hline
  {\scriptsize\ttfamily\textbackslash UTIDatosTutor} & Registra el tutor o tutora y su género gramatical. \\
  \hline
  {\scriptsize\ttfamily\textbackslash UTIPortada} & Registra facultad, carrera, tema, modalidad, ciudad, año y genera la portada. \\
  \hline
\end{longtable}

\section{Líneas de investigación}

El catálogo se encuentra en \texttt{uti\_tfm.sty}. La selección se realiza con dos índices. El primer índice corresponde a la línea y el segundo a la sublínea. El valor \texttt{\{0\}\{0\}} mantiene la selección pendiente.

La tabla de Área de Estudio consume automáticamente la línea y sublínea seleccionadas. El campo se fija institucionalmente como ``Tecnologías de la información y la comunicación (TIC)''.

\section{Área de estudio}

La tabla institucional incluye dominio, línea, sublínea, campo, área, aspectos, objeto y periodo de estudio. Los campos no parametrizados se editan en \texttt{Capítulos/02\_Metodología.tex}.

\chapter{Elementos técnicos y numeración}

\section{Figuras, tablas e imágenes}

Las figuras, tablas, imágenes y ecuaciones usan numeración arábiga continua; no reinician al cambiar de capítulo. Las leyendas tienen etiqueta en negrita, separador de dos puntos y texto justificado.

La plantilla distingue:

\begin{itemize}
  \item \texttt{figure}: diagramas, gráficos, arquitecturas o esquemas;
  \item \texttt{table}: datos organizados en filas y columnas;
  \item \texttt{imagen}: fotografías, capturas de pantalla o evidencias visuales;
  \item \texttt{equation}: expresiones matemáticas numeradas.
\end{itemize}

Los comandos \texttt{\textbackslash elaboradopor} y \texttt{\textbackslash fuente} registran autoría y procedencia. Toda pieza debe poseer una etiqueta única y ser mencionada mediante \texttt{\textbackslash ref} antes o después de su aparición.

\section{Ecuaciones y anexos}

Las ecuaciones se añaden manualmente a su índice con el comando \texttt{addeqtolist}. Los anexos se crean con \texttt{Anexo}, lo cual genera la letra correspondiente y actualiza el índice correspondiente.

\chapter{Referencias y compilación}

\section{Sistema bibliográfico}

La versión actual utiliza BibTeX, no biber. El estilo bibliográfico es \texttt{IEEEtran} y la base utilizada es \texttt{referencias.bib}.

Las citas se insertan mediante el comando \texttt{cite}. Todas las entradas de \texttt{referencias.bib} deben poseer autor, título, año y los campos exigidos por su tipo documental.

\section{Secuencia de compilación}

Para XeLaTeX:

\begin{lstlisting}
xelatex main.tex
bibtex main
xelatex main.tex
xelatex main.tex
\end{lstlisting}

Para pdfLaTeX:

\begin{lstlisting}
pdflatex main.tex
bibtex main
pdflatex main.tex
pdflatex main.tex
\end{lstlisting}

La alternativa recomendada para automatizar el proceso es:

\begin{lstlisting}
latexmk -xelatex main.tex
% o, si Century Gothic no esta instalada:
latexmk -pdf main.tex
\end{lstlisting}

\chapter{Gobierno, control de cambios y mantenimiento}

\section{Elementos controlados por el Comité}

Requieren aprobación previa:

\begin{itemize}
  \item márgenes, tamaño de papel y tipografías;
  \item contenido y orden de páginas legales;
  \item estructura obligatoria de capítulos;
  \item formato de portada, resumen y abstract;
  \item sistema de citación y estilo bibliográfico;
  \item catálogo de líneas y sublíneas de investigación;
  \item reglas de numeración e índices;
  \item denominación de títulos profesionales y carrera.
\end{itemize}

\section{Elementos editables por el estudiante}

El estudiante puede modificar metadatos, contenido académico, recursos visuales, bibliografía, anexos y campos descriptivos del Área de Estudio. No debe alterar \texttt{uti\_tfm.sty} ni eliminar páginas obligatorias sin autorización.

\section{Procedimiento de cambio}

\begin{enumerate}
  \item Registrar la necesidad y su fundamento académico o normativo.
  \item Identificar los archivos y documentos afectados.
  \item Preparar una versión de prueba sin alterar la versión aprobada.
  \item Compilar casos de uno y dos autores.
  \item revisar portada, preliminares, índices, capítulos, referencias y anexos.
  \item Presentar evidencia PDF al Comité Curricular.
  \item Aprobar, asignar nueva versión y actualizar ambos manuales.
\end{enumerate}

\chapter{Criterios de aceptación}

\section{Pruebas funcionales mínimas}

\begin{longtable}{|p{4.1cm}|p{5.2cm}|p{3.3cm}|}
  \hline
  \textbf{Prueba} & \textbf{Criterio de aceptación} & \textbf{Evidencia} \\
  \hline
  Compilación limpia & El documento genera PDF sin errores que detengan el proceso. & Registro de compilación. \\
  \hline
  Márgenes & Se conservan 3/3/4/3 cm en superior, inferior, izquierdo y derecho. & Inspección del PDF. \\
  \hline
  Un autor & Pronombres, firmas y páginas preliminares aparecen en singular. & Caso de prueba. \\
  \hline
  Dos autores & Pronombres, firmas, dedicatorias y agradecimientos se adaptan al plural. & Caso de prueba. \\
  \hline
  Tutor o tutora & La etiqueta y el encabezado cambian según el género configurado. & Caso de prueba. \\
  \hline
  Línea de investigación & Un par válido carga la línea y sublínea correctas en Área de Estudio. & Tabla generada. \\
  \hline
  Índices & Contenidos, tablas, figuras, imágenes, ecuaciones y anexos se actualizan. & Páginas de índices. \\
  \hline
  Referencias & Las citas se numeran en orden y la bibliografía adopta IEEE. & Sección Referencias. \\
  \hline
  Abstract & Presenta datos en inglés y usa el título inglés configurado. & Página de Abstract. \\
  \hline
\end{longtable}

\section{Lista de aprobación curricular}

\begin{itemize}
  \item[$\square$] La portada cumple la identidad y jerarquía institucional.
  \item[$\square$] Los textos legales han sido revisados por la instancia competente.
  \item[$\square$] La estructura académica corresponde al perfil de la carrera.
  \item[$\square$] Las líneas de investigación coinciden con el catálogo vigente.
  \item[$\square$] El formato IEEE ha sido ratificado para citas y referencias.
  \item[$\square$] Los casos de uno y dos autores fueron validados.
  \item[$\square$] Los dos motores de compilación previstos fueron evaluados.
  \item[$\square$] El Manual de Uso es consistente con esta versión.
\end{itemize}

\chapter*{Registro de aprobación}
\addcontentsline{toc}{chapter}{Registro de aprobación}

\vspace{1cm}

\begin{longtable}{|p{4cm}|p{4cm}|p{4cm}|}
  \hline
  \textbf{Rol} & \textbf{Nombre y firma} & \textbf{Fecha} \\
  \hline
  Coordinación de carrera & & \\
  \hline
  Coordinación de titulación & & \\
  \hline
  Presidencia del Comité Curricular & & \\
  \hline
  Secretaría del Comité Curricular & & \\
  \hline
\end{longtable}

\vfill
\noindent\textbf{Observaciones de aprobación:}

\vspace{3cm}
\hrule

\end{document}